% General log canonical extension, separate from the checked Gorenstein fallback.
% Root, C and D independent actual-file audits passed.

\begin{lemma}[Proved: exceptional configurations for arbitrary lc surfaces]
\label{D:lc-exceptional-configurations}
Let $(S,s)$ be a normal log canonical surface singularity over
an algebraically closed field, in arbitrary characteristic.
On the minimal regular resolution its reduced exceptional
divisor has one of the following forms:
\begin{enumerate}
\item a simple normal crossings tree of smooth rational curves;
\item one smooth elliptic curve, with discrepancy $-1$;
\item a rational nodal cycle, with discrepancy $-1$ on every
component, allowing one irreducible nodal rational curve.
\end{enumerate}
The exceptional divisor is empty if the point is regular.
No Gorenstein hypothesis or restriction on the Cartier
index is imposed.  In the tree case we do not assert that
the singularity is rational or log terminal, or that its
discrepancies are zero.
\end{lemma}
\begin{proof}
Let $f:Z\to S$ be the minimal regular resolution and write
\[
 K_Z=f^*K_S+\sum_i a_i E_i.
\]
The exceptional fibre is connected and its intersection
matrix is negative definite.  The surface facts and exact
sources used here are the ones in
Lemma~\ref{D:Gorenstein-lc-configurations}:
\cite[section 10.1, Theorem 10.1]{Kollar2013D} and
\cite[tags 0BGP, 0C2M and 0C5R]{StacksDSurfaceResolution}.
Adjunction and minimality give
\[
 K_Z\cdot E_i=2p_a(E_i)-2-E_i^2\geq0.
\]
The positive-part argument for the negative definite
intersection matrix therefore gives $a_i\leq0$, exactly
as in the proof of that lemma.  Log canonicity gives
$a_i\geq-1$.  We use only
\begin{equation}\label{D:eq-lc-discrepancy-interval}
                  -1\leq a_i\leq0;
\end{equation}
the coefficients need not be integral.

We first record a useful consequence of a zero discrepancy.
If $a_j=0$, then
\[
 0\leq K_Z\cdot E_j
    =\sum_{i\ne j}a_i(E_i\cdot E_j)\leq0.
\]
Hence every component adjacent to $E_j$ also has
discrepancy zero.  In particular a zero-discrepancy
component cannot meet a component of negative discrepancy.

Suppose that some $E_i$ has arithmetic genus at least one.
From \eqref{D:eq-lc-discrepancy-interval} we obtain
\[
 -E_i^2\leq K_Z\cdot E_i
       =\sum_j a_j(E_j\cdot E_i)
       \leq a_iE_i^2\leq -E_i^2.
\]
All terms are equal.  Thus $p_a(E_i)=1$, $a_i=-1$,
and every component meeting $E_i$ would have discrepancy
zero.  The preceding observation excludes all such
neighbors.  By connectedness, $E_i$ is the whole fibre.
We will determine its possible singularities below.

We may now suppose that every component has arithmetic
genus zero, and hence is a smooth rational curve.
Write
\[
 b_i=-E_i^2\geq2,\qquad c_i=1+a_i\in[0,1].
\]
The lower bound on $b_i$ is minimality.
Form the intersection multigraph, placing
$E_i\cdot E_j$ edges between distinct vertices $i,j$.
This counts intersection multiplicities, before any
normal-crossings property has been established.
If this multigraph contains a cycle, select one with
distinct vertices, allowing a two-vertex cycle consisting
of two selected parallel edges.
Let $I$ be its vertex set and $s_{ij}$ the multiplicity
of its selected edges.  Thus
$\sum_{j\in I}s_{ij}=2$ for every $i\in I$.
For internal extra edges put
$r_{ij}=E_i\cdot E_j-s_{ij}\geq0$.
For an external vertex $j\notin I$ put
$m_{ij}=E_i\cdot E_j$.

Adjunction, or equivalently intersecting
$K_Z+\sum_iE_i=f^*K_S+\sum_i c_i E_i$
with $E_i$, gives, for $i\in I$,
\begin{equation}\label{D:eq-cycle-extra-edges}
 b_i c_i-\sum_{j\in I}s_{ij}c_j
 =-\sum_{j\in I}r_{ij}(1-c_j)
  -\sum_{j\notin I}m_{ij}(1-c_j).
\end{equation}
Summing over the selected cycle yields
\begin{equation}\label{D:eq-cycle-sum}
 \sum_{i\in I}(b_i-2)c_i
 =-\sum_{\{i,j\}\subset I}r_{ij}(2-c_i-c_j)
  -\sum_{\substack{i\in I\\j\notin I}}m_{ij}(1-c_j).
\end{equation}
Every term on the left is nonnegative and every term
on the right is nonpositive.  All are therefore zero.
In particular, each extra internal edge would have
both endpoints $c_i=c_j=1$, and every external
neighbor of the selected cycle would have $c_j=1$.

Negative definiteness guarantees some $b_i>2$ on the
selected cycle.  Otherwise the vector equal to one
on $I$ and zero elsewhere would have square
\[
 -2|I|+2|I|+2\sum_{\{i,j\}\subset I}r_{ij}\geq0,
\]
which is impossible.  Thus $c_i=0$ at some vertex.
All the right-hand terms in
\eqref{D:eq-cycle-extra-edges} have already vanished,
so the equations reduce to
$b_i c_i=\sum_j s_{ij}c_j$.
Nonnegativity propagates this zero around the selected
cycle.  Hence $c_i=0$ for every $i\in I$, or
$a_i=-1$ there.
An extra internal edge is now impossible, since its
endpoints would have to be one.  An external neighbor
would have $a_j=0$, impossible next to $a_i=-1$ by
the zero-discrepancy observation.  Connectedness
therefore shows that the whole exceptional divisor
is this cycle, with all discrepancies $-1$.
This argument covers intersection multiplicities
larger than two; none were discarded as if they
were single edges.

We check nodality in both cases in which all
discrepancies have been shown to be $-1$.
For the reduced exceptional curve $E$ we have
$K_Z+E=f^*K_S$, so $(Z,E)$ is a log canonical pair.
At a point of multiplicity at least three, the
point blowup would have discrepancy at most $-2$.
At a multiplicity-two point with one repeated
tangent, blow up once.  The strict transform
meets the new discrepancy-$-1$ exceptional curve
at one point with intersection multiplicity two.
If it is singular there, their total multiplicity
is at least three.  If it is smooth, it is tangent
to that exceptional curve, and the next blowup
creates a point with total reduced-boundary
multiplicity at least three.  A subsequent blowup
again gives discrepancy at most $-2$.
Thus every singularity of $E$ is an ordinary node,
in every characteristic.
An integral nodal curve of arithmetic genus one
is either smooth elliptic or a rational curve
with one node.  In the rational multigraph cycle
case, nodality gives exactly a rational nodal
cycle, including two components meeting at two
distinct nodes.

Finally, if the rational intersection multigraph
has no cycle, it is a tree.  There can be no
intersection multiplicity greater than one, since
this would give a two-edge cycle, and no triple
point, since this would give a three-cycle.
All components are smooth rational and all
intersections are transverse, so it is a simple
normal crossings rational tree.  This proves all
the alternatives.
\end{proof}

\begin{theorem}[Proved: tame ramification above an arbitrary lc cusp]
\label{D:lc-cusp-tame}
Let $(S,s)$ be a normal log canonical surface singularity
over an algebraically closed field.  Suppose that the
minimal exceptional fibre is a rational nodal cycle.
Every finite normal separable cover of $S$ that is
\'etale on the punctured germ is tame at every component
of that cycle, and at the components obtained by
subdividing its nodes.  No Gorenstein or index
restriction is required.
\end{theorem}
\begin{proof}
We indicate explicitly which argument replaces the
Gorenstein input in Theorem~\ref{D:cusp-tame}, and
repeat the decisive comparison.
The target discrepancies on the cycle are $-1$ by
Lemma~\ref{D:lc-exceptional-configurations}.
Pass to a Galois closure of the punctured cover
and normalize.  The canonical-divisor equality
$K_{S'}=\pi^*K_S$ holds as an equality of rational
Cartier divisors, because the cover is \'etale
in codimension one.  It makes $S'$
$\mathbb Q$-Gorenstein.  For every divisorial
valuation extension $F/E$, the determinant
calculation of Proposition~\ref{D:lc-elliptic-tame}
gives
\[
 a(F,S')=e(F/E)a(E,S)+\delta_F,
 \qquad \delta_F\geq e(F/E)-1.
\]
Consequently $S'$ is log canonical.  No
integrality of its discrepancies is asserted.
We may apply the new configuration lemma to
its minimal resolution.

Work in the strict henselian local germ, choosing
a connected covering component and its decomposition
group.  The finiteness, normality and residue-field
comparisons are exactly those verified in
Theorem~\ref{D:cusp-tame}.  Let $Y\to S$ be a
regular resolution with smooth-component cycle
$E$, after one node blowup if needed.
Proper henselian comparison
\cite[Lemma 58.9.1, tag 0A48]{StacksBCuspLocal}
and normalization of the nodal curve give
\[
 \pi_1(Y)\simeq\pi_1(E)
       \simeq\widehat{\mathbb Z}.
\]
The image of the finite-index subgroup
$\pi_1(S'\setminus\{s'\})$ in this quotient
is again finite-index and nonzero.

Resolve the graph from a minimal regular resolution
of $S'$ to $Y$.  Let $Z\to S'$ be the resulting
regular resolution and let $E'$ be its reduced
exceptional fibre.  Functorial proper henselian
comparison shows that $\pi_1(E')\to\pi_1(E)$
has that nonzero image.  Point-blowup factorization
shows that $E'$ is obtained from its minimal core
by rational tree attachments and cycle-node
subdivisions.
If the covering configuration is a rational tree
or a regular point, then $\pi_1(E')=1$.
If it is elliptic, its irreducible elliptic core
maps into one rational component of $E$, so
its fundamental-group image is again zero.
Both are impossible.  Thus the covering
configuration is a rational cycle, and the
new classification gives discrepancy $-1$
on every component of its cyclic core.

The source cyclic core must dominate every
component of $E$: otherwise its connected image
lies in a proper rational chain of $E$, which
has trivial fundamental group.  This would
contradict the nonzero image just established.
For each component $E_i$, we have therefore
found an extension $F_i$ of its valuation with
$a(F_i,S')=-1$.
The Galois group acts transitively on the
valuation extensions of $E_i$ in the normalized
cover of $Y$.  Discrepancy is invariant under
these deck transformations.  Every extension
$F/E_i$ thus has discrepancy $-1$, and
\[
 -1=-e(F/E_i)+\delta_F.
\]
Equality $\delta_F=e(F/E_i)-1$ says precisely
that the extension is tame, including residue
separability.  The strict henselian comparison
preserves these indices and residue extensions,
so the conclusion holds for the algebraic germ.
The same proof begins with any subdivision of
the target cycle and proves the final assertion.
\end{proof}

\begin{corollary}[Proved: all valuations above an lc elliptic or cyclic core]
\label{D:lc-core-valuative-tame}
If the minimal exceptional fibre of a normal
log canonical surface singularity over an
algebraically closed field is elliptic or cyclic,
then every finite normal separable cover that
is \'etale on the punctured germ is tame at every
divisorial valuation centred at the singular point.
No restriction on the Cartier index is imposed.
\end{corollary}
\begin{proof}
After blowing up a possible self-node, the
minimal core is a simple normal crossings
divisor.  Proposition~\ref{D:lc-elliptic-tame}
and Theorem~\ref{D:lc-cusp-tame} give tameness
on its components.  The Kummer persistence
argument of Lemma~\ref{D:Gorenstein-lc-valuative-tame}
uses only this regular normal-crossings model
and its componentwise tameness.  Applying it
here proves tameness on every further point
blowup, and hence at every specified exceptional
divisorial valuation.  The cited argument is
\cite[XIII, Proposition 5.2 and Corollary 5.3,
pp.~428--430]{SGA1}; it uses no Cartier-index
condition after componentwise tameness is known.
\end{proof}
